\section{Introduction}

Let \(Q\subset\R^2\) be a nondegenerate Euclidean triangle. Its unit-speed
billiard flow \(\Phi_t\) moves in straight lines in the interior and reflects
specularly at the open sides. Trajectories meeting a vertex in either time
direction are discarded; they form a null set. The invariant probability
measure on position and direction is normalized Liouville measure
\[
 d\mu(x,\theta)=\frac{dA(x)\,d\theta}{2\pi\operatorname{Area}(Q)}.
\]
We say that the flow is \emph{weakly mixing} if its product with itself,
\((x,y)\mapsto(\Phi_t x,\Phi_t y)\), is ergodic for
\(\mu\otimes\mu\). Thus every measurable set invariant under all product
time maps has measure zero or one.

\begin{theorem}\label{thm:main}
The unit-speed billiard flow in every nondegenerate Euclidean triangle with
at least one angle irrational relative to \(\pi\) is weakly mixing for
normalized Liouville measure.
\end{theorem}

In particular, for triangles of area one with angles
\(\alpha,\beta,\pi-\alpha-\beta\), the conclusion holds for
Lebesgue-almost every point of
\[
 D=\{(\alpha,\beta)\in\R^2:\alpha>0,\ \beta>0,\ \alpha+\beta<\pi\}:
\]
the excluded pairs with both \(\alpha/\pi\) and \(\beta/\pi\) rational
form a countable set. The irrationality assumption is sharp for the full
phase space. In a rational triangle, side reflections generate a finite
group on directions. A nonconstant function on the circle invariant under
this finite group is then an invariant function for the billiard flow, so
the full flow is not ergodic.

\paragraph{Historical context.}
Polygonal billiards connect elementary reflection geometry with dynamics
on flat surfaces. The unfolding construction of Zemlyakov and Katok
\cite{ZK} replaces reflections by straight continuation on a surface.
For rational polygons, it produces a compact translation surface and
separates the full phase space into invariant directional surfaces.
Kerckhoff, Masur, and Smillie proved unique ergodicity in almost every
direction \cite{KMS}. They also obtained ergodicity of the full billiard
flow for a dense \(G_\delta\) family of polygons. Vorobets gave a
quantitative rational-approximation criterion and explicit ergodic
examples, including irrational right triangles \cite{Vorobets}.
These results establish ergodicity under genericity or sufficiently rapid
approximation assumptions rather than for each irrational triangle.

Chaika and Forni proved that the full billiard flow is weakly mixing for a
dense \(G_\delta\) set of polygons with any fixed number of sides at least
three \cite{CF}. Their argument uses the absence of common nontrivial
eigenvalues for directional flows in almost every pair of directions on a
translation surface. The Baire-category conclusion does not assert
Lebesgue-almost-everywhere weak mixing in the space of triangle angles.
For rational billiards, recent work of Arana-Herrera, Chaika, and Forni
characterizes weak mixing in almost every direction \cite{ACF}; this is
a statement on directional invariant surfaces, rather than on the full
Liouville phase space considered here.

Numerical investigations have long suggested weak mixing in irrational
right triangles and mixing in triangles with all angles irrational
\cite{ACG,CP}. Other computations have suggested nonergodicity in some
irrational right or symmetric triangles \cite{WCP,ZSBJ}. These conflicting
finite-time observations underscore the difficulty of identifying
measurable invariants from trajectories with slow statistical convergence.
Theorem~\ref{thm:main} treats every irrational triangle without an
approximation or symmetry restriction.

\paragraph{Analytic contribution.}
The regularity of an eigenfunction is the central issue. Glue two oppositely
oriented copies of \(Q\) along their sides. The resulting flat double is
regular across the open sides and has cone points at the vertices. On its
unit tangent bundle, let \(X\) denote differentiation along a geodesic and
\(Y\) differentiation in the perpendicular parallel direction. Locally
these vector fields commute. Forni and Moll developed angular Fourier and
Cauchy--Riemann methods in this setting \cite[Sections~3.2--3.5]{FM}.
Their rigidity results distinguish invariance under both horizontal
directions from geodesic invariance under an additional horizontal Sobolev
hypothesis \cite[Theorem~4.1 and Corollary~4.2]{FM}.
A measurable flow eigenfunction supplies only \(Xf=i\lambda f\) at the
outset; the missing transverse derivative cannot be assumed.

The companion manuscript \cite[Sections~2--6]{Companion} removes this
regularity obstacle for invariant functions, proving ergodicity of
irrational triangular billiards. We extend its coefficient-energy and
long-average method to every real spectral parameter \(\lambda\).
Summing the coefficient-energy identities over consecutive angular indices
makes the terms containing \(\lambda\) telescope to a single boundary term.
A separate pairing with the two boundary modes produced by differentiating
an angular Fourier tail gives the sharp bound
\[
 \norm{Xf}_2^2+\norm{Yf}_2^2\le\lambda^2\norm{f}_2^2.
\]
Since \(Xf=i\lambda f\), it follows that \(Yf=0\).
The use of a spectral parameter in Cauchy--Riemann energy estimates has a
related precedent in the twisted cohomological equations studied by Forni
\cite[Section~3]{Forni2022}. Here the argument starts with bounded measurable
eigenfunctions and proves the needed derivative bounds. We give a
self-contained proof of this analytic extension below.

\paragraph{Proof strategy.}
Expand a bounded eigenfunction on the double in angular Fourier modes.
Spatial smoothing away from the cone points and cutoffs near the tips
first give a uniform bound on the gradient of each coefficient. These
individual bounds do not
yet give a square-summable family of gradients. To obtain the missing
summability, take the sum of all modes above a fixed index. Its derivative
under \(X-i\lambda\) consists of two endpoint terms.

We test this equation using long time averages supported on orbit segments
that stay clear of the cone points. There are two approximation parameters.
At a fixed bound on the discarded measure, the smoothing limit gives a
bounded solution of the exact local eigenfunction equation and passes the
endpoint pairing to it.
The individual-coefficient estimate then supplies a new derivative bound
independent of the discarded measure. A second limit passes the pairing
to the original eigenfunction. Summing the resulting endpoint inequalities
proves \(Yf=0\). This order of limits is what permits the initial
eigenfunction to be merely measurable.

The remaining argument is geometric. On the triangle interior,
\(Xf=i\lambda f\) and \(Yf=0\) give a plane wave with a measurable
direction-dependent amplitude. Reflection in the two sides through an
irrational vertex makes this amplitude constant. Reflection in the third
side excludes every nonzero frequency. The standard spectral criterion
then gives weak mixing.

Section~\ref{sec:geometry} sets up the double and its Fourier coefficients.
Section~\ref{sec:energy} proves the coefficient and smoothing estimates.
Section~\ref{sec:rigidity} establishes transverse rigidity by the two
successive limits. Section~\ref{sec:conclusion} applies side reflections
and proves ergodicity of the product flow.
